\subsection{Geometrische Definition}
\label{subsec:vector_algebra_scalar_product_definition}

\begin{DefinitionBox}{Skalarprodukt}{vector_algebra_scalar_product}
  \begin{minipage}[c]{0.65\linewidth}
    Das Skalarprodukt \EnglishTerm{dot product} zweier Vektoren $\vec a,\vec b\in\R^n$
    mit einem Zwischenwinkel $\varphi$ ist definiert durch
    \[
      \vec a\cdot\vec b
      =|\vec a|\,|\vec b|\cos(\varphi)\in\R,
    \]
    Das Ergebnis eines Skalarprodukts ist immer ein Skalar!
  \end{minipage}\hfill
  \begin{minipage}[c]{0.33\linewidth}
    \centering
    \CachedTikzFigure{content/resources/chapters/vector_algebra/tikz/dot_product_geometry.tex}
  \end{minipage}
\end{DefinitionBox}

\begin{ExampleBox}{Skalarprodukt aus Längen und Winkel}{vector_algebra_geometric_dot_product_example}
  \noindent
  \begin{minipage}[c]{0.65\linewidth}
    Für zwei Vektoren mit den Längen $|\vec a|=2$ und 
    $|\vec b|=3$ sowie einem Zwischenwinkel $\varphi=30^\circ$
    ist das Skalarprodukt
    \[
      \vec a\cdot\vec b
      =\Cloze{8cm}{
        2\cdot3\cdot\cos(30^\circ) = 2\cdot3\cdot\frac{\sqrt3}{2} = 3\sqrt3
        }.
    \]
  \end{minipage}\hfill
  \begin{minipage}[c]{0.33\linewidth}
    \centering
    \CachedTikzFigure{content/resources/chapters/vector_algebra/tikz/geometric_dot_product_example.tex}
  \end{minipage}
\end{ExampleBox}

Aus Definition~\ref{def:vector_algebra_scalar_product} lassen sich unmittelbar
sehr wichtige Folgerungen ableiten.
\begin{KeyPointBox}{Skalarprodukt senkrechter Vektoren}{vector_algebra_dot_product_orthogonal_vectors}
  \begin{minipage}[c]{0.68\linewidth}
    Für senkrechte Vektoren ist der Zwischenwinkel $\varphi=90^\circ$.
    Da $\cos(90^\circ)=0$ folgt für das Skalarprodukt
    \[
      \Cloze{9.5cm}{%
        \vec a\perp\vec b
        \quad \Longleftrightarrow \quad
        \vec a\cdot\vec b=0.}
    \]
    Zwei Vektoren sind genau dann \emph{orthogonal} \EnglishTerm{orthogonal},
    wenn ihr Skalarprodukt null ist.
  \end{minipage}\hfill
  \begin{minipage}[c]{0.27\linewidth}
    \centering
    \CachedTikzFigure{content/resources/chapters/vector_algebra/tikz/orthogonal_vectors_dot_product.tex}
  \end{minipage}
\end{KeyPointBox}

\begin{KeyPointBox}{Skalarprodukt paralleler Vektoren}{vector_algebra_dot_product_parallel_vectors}
  \begin{minipage}[c]{0.68\linewidth}
    Für gleichgerichtete parallele Vektoren ist der Zwischenwinkel
    $\varphi=0^\circ$. Aus
    Definition~\ref{def:vector_algebra_scalar_product} folgt damit
    \[
      \Cloze{9.5cm}{%
        \vec a\parallel\vec b
        \quad \Longleftrightarrow \quad
        \vec a\cdot\vec b=|\vec a|\,|\vec b|.}
    \]
    Das Skalarprodukt paralleler Vektoren entspricht dem Produkt der Beträge. \\
  \end{minipage}\hfill
  \begin{minipage}[c]{0.27\linewidth}
    \centering
    \CachedTikzFigure{content/resources/chapters/vector_algebra/tikz/parallel_vectors_dot_product.tex}
  \end{minipage}
    \textbf{Spezialfall:} Für $\vec b=\vec a$ ergibt sich
    \[
      \Cloze{9.5cm}{%
        \vec a\cdot\vec a=|\vec a|^2
        \quad\Longleftrightarrow\quad
        |\vec a|=\sqrt{\vec a\cdot\vec a}.}
    \]
\end{KeyPointBox}

Zusammenfassend bestimmt der Zwischenwinkel das Vorzeichen des Skalarprodukts.

\begin{KeyPointBox}{Vorzeichen und eingeschlossener Winkel}{vector_algebra_dot_product_sign}
  \begin{center}
    \CachedTikzFigure{content/resources/chapters/vector_algebra/tikz/dot_product_signs.tex}
  \end{center}
\end{KeyPointBox}

\begin{ExampleBox}{Skalarprodukte kanonischer Einheitsvektoren}{vector_algebra_canonical_unit_vectors_dot_products}
  Wir unterscheiden zwischen Skalarprodukten gleicher und unterschiedlicher kanonischer Einheitsvektoren:
  \begin{enumerate}[label=\alph*),leftmargin=*,itemsep=0.8em]
    \item \textbf{Gleiche Einheitsvektoren ($i=j$):}
      Da $\vec e_i$ zu sich selbst parallel ist und den Betrag
      $|\vec e_i|=1$ besitzt, folgt
      \[
        \vec e_i\cdot\vec e_i
          =|\vec e_i|\,|\vec e_i|
          =1\cdot1
          =\Cloze{1cm}{1}.
      \]
    \item \textbf{Unterschiedliche Einheitsvektoren ($i\ne j$):}
      Die Vektoren stehen senkrecht aufeinander.
      \[
        \vec e_i\perp\vec e_j
          \quad\Longrightarrow\quad
        \vec e_i\cdot\vec e_j=\Cloze{1cm}{0}.
      \]
  \end{enumerate}

  Zusammengefasst gilt
  \[
    \vec e_i\cdot\vec e_j
      =\begin{cases}
          \Cloze{0.8cm}{1},&i=j,\\
          \Cloze{0.8cm}{0},&i\ne j.
        \end{cases}
  \]
\end{ExampleBox}
